% ---------------------------------------------------------------------------
% Document macros.
%
% These belong in a preamble, but the authoritative main.tex preamble is fixed
% and must not be edited, so they live here and are \input'ed as the first
% thing inside the document body. \newcommand is legal in the body, and both
% macros are expanded only at use sites that come after this \input. If the
% authoritative preamble is ever reopened, move these two definitions into it
% verbatim and delete this file plus its \input line.
% ---------------------------------------------------------------------------

% System name. The original definition was \newcommand{\sysname}{Tandem\xspace}
% but `xspace' is NOT loaded by the authoritative preamble. Rather than add a
% package, the macro is plain and every call site that is followed by a space
% or a newline was rewritten to \sysname{} (mechanically, 20 sites across three
% section files). Call sites followed directly by punctuation or a closing
% brace need no {} and were left alone.
\newcommand{\sysname}{Tandem}

% ---------------------------------------------------------------------------
% FLOAT PLACEMENT. This paper is figure-dense: 14 floats against roughly six
% pages of text. LaTeX's stock parameters allow at most 2 floats at the top of
% a page and 3 per page total, and refuse any page that is more than 70% float
% -- limits tuned for a paper with far fewer figures. Under them the float
% queue cannot drain as fast as figures are declared, and the overflow is
% flushed at the end of the document, which is why earlier drafts ended in a
% block of figures.
%
% Raising the limits lets the queue keep up. These are placement parameters
% only: nothing about the text changes, and a figure still cannot appear
% before the point where it is DECLARED (which is why the section files now
% declare each float at the head of its subsection rather than after the
% paragraph that cites it).
%
% If figures still collect at the end after this, the queue is genuinely
% oversubscribed and the remedy is fewer or merged figures, not more tuning.
\setcounter{topnumber}{3}            % floats at the top of a column  (was 2)
\setcounter{totalnumber}{4}          % floats per column              (was 3)
\setcounter{dbltopnumber}{2}         % full-width floats at page top  (was 2)
\renewcommand{\topfraction}{0.92}    % of a column that may be top float
\renewcommand{\dbltopfraction}{0.92} % same, for full-width floats
\renewcommand{\bottomfraction}{0.5}
\renewcommand{\textfraction}{0.06}   % minimum text on a mixed page   (was .2)
\renewcommand{\floatpagefraction}{0.7}
\renewcommand{\dblfloatpagefraction}{0.7}

% Gap between a float (including its caption) and the body text sharing the
% column, cut to 70% of whatever IEEEtran set. Written as a self-multiplication
% rather than an absolute length so it scales IEEEtran's own values -- which
% are expressed in \baselineskip and carry stretch and shrink components -- and
% keeps that elasticity instead of pinning the glue rigid. Pinning it would
% cost more space than it saves, because the output routine loses the give it
% uses to balance columns.
%
%   \textfloatsep     float at a column top or bottom <-> body text
%   \floatsep         between two floats stacked in one column
%   \dbltextfloatsep  same as \textfloatsep, for full-width floats
%   \dblfloatsep      between two stacked full-width floats
%   \intextsep        around an [h] float (none are used here; set for safety)
%
% \abovecaptionskip is deliberately untouched: that is the gap between the
% graphic and its own caption, and shrinking it detaches neither, it just
% crowds them. Add \setlength{\abovecaptionskip}{0.7\abovecaptionskip} here if
% more is needed after seeing a build.
\setlength{\textfloatsep}{0.7\textfloatsep}
\setlength{\floatsep}{0.7\floatsep}
\setlength{\dbltextfloatsep}{0.7\dbltextfloatsep}
\setlength{\dblfloatsep}{0.7\dblfloatsep}
\setlength{\intextsep}{0.7\intextsep}

% ---------------------------------------------------------------------------
% VERTICAL GLUE. In two-column mode LaTeX makes both columns the same height by
% stretching whatever glue it can find -- between paragraphs, around section
% headings, around displays. On a page with several floats there is little text
% left to absorb that, so the stretch lands in a few large gaps: the "ton of
% inter-paragraph space" around short subsections.
%
% \raggedbottom tells the output routine not to do that. Columns are set at
% their natural height and end where they end, which is standard for IEEE
% conference papers and is the difference between a short subsection reading as
% short and reading as a hole in the page.
\raggedbottom

% Space around a displayed equation. IEEEtran's defaults assume a display is a
% visual break; Eq. 1 here is read as part of the sentence around it, so it
% does not need the full separation.
\setlength{\abovedisplayskip}{4pt plus 2pt minus 2pt}
\setlength{\belowdisplayskip}{4pt plus 2pt minus 2pt}
\setlength{\abovedisplayshortskip}{2pt plus 2pt}
\setlength{\belowdisplayshortskip}{3pt plus 2pt minus 2pt}
% ---------------------------------------------------------------------------

% Placeholder that used to stand in for an unrendered plot. #1 optional width,
% #2 box height, #3 the plot specification.
%
% AS OF THIS REVISION IT HAS NO CALL SITES. All 14 figures carry a real
%     \includegraphics[width=\columnwidth]{figures/NAME}
% (0.95\textwidth for fig:budget-sweep, the only figure*), resolved relative to
% main.tex -- the authoritative preamble sets no \graphicspath and none is
% needed. Every figure is generated by scripts/figures/make_figures.py; none is
% author-supplied. The plot specification each placeholder carried survives as
% a LaTeX comment above its \includegraphics, since that is the only written
% record of what the figure is meant to show.
%
% Kept only so an unconverted draft still compiles. Safe to delete.
\newcommand{\figspec}[3][0.95\columnwidth]{%
  \fbox{\parbox[c][#2][c]{#1}{\centering\footnotesize\itshape #3}}}
